% Generated draft: requires mathematical and source review.
\documentclass[11pt]{article}
\usepackage[margin=25mm]{geometry}
\usepackage{fontspec}
\setmainfont{Latin Modern Roman}
\usepackage{amsmath,amssymb,mathtools,unicode-math}
\setmathfont{Latin Modern Math}
\usepackage{enumitem,needspace,xurl,hyperref}
\setlength{\parindent}{0pt}
\setlength{\parskip}{5pt}
\setlength{\emergencystretch}{3em}
\newcommand{\N}{\mathbb{N}}
\newcommand{\Z}{\mathbb{Z}}
\newcommand{\Q}{\mathbb{Q}}
\newcommand{\R}{\mathbb{R}}
\newcommand{\C}{\mathbb{C}}
\newcommand{\F}{\mathbb{F}}
\newcommand{\A}{\mathbb{A}}
\newcommand{\PP}{\mathbb{P}}
\newcommand{\E}{\mathbb{E}}
\newcommand{\eps}{\varepsilon}
\newcommand{\dd}{\,\mathrm{d}}
\begin{document}
\section*{UnsolvedMath ID: 30000864}
\textbf{Source problem code:} catalogue entry
\subsection{Definitions and mathematical statement}
Let $m\geq1$ and let $k>2m$ be an integer, with $l=k-2m$. Let $f,g$ be the source's Siegel modular forms of genus $2m$ and weights $k,l$, respectively. To attach Satake parameters, they must be understood as nonzero simultaneous Hecke eigenforms at the relevant unramified primes. For each such prime $p$, write their parameter tuples as
\[
(\alpha_0(p),\alpha_1(p),\ldots,\alpha_{2m}(p)),\qquad
(\beta_0(p),\beta_1(p),\ldots,\beta_{2m}(p)).
\]
Satake tuples are understood with a common normalization and up to the relevant Weyl-group action.

The conjecture asks for a single Siegel Hecke eigenform $F$ of genus $4m$ and weight $k$ whose Satake representatives can be chosen at every intended common unramified prime so that
\[
\begin{aligned}
\gamma_0(p)&=\alpha_0(p)\beta_0(p),\\
\gamma_i(p)&=\alpha_i(p)\qquad(1\leq i\leq2m),\\
\gamma_{2m+j}(p)&=\beta_j(p)\qquad(1\leq j\leq2m).
\end{aligned}
\]
The weights and both parameter ranges are exactly those in the source. The conclusion is a global form simultaneously realizing the local tuples, not a separately chosen object for each prime. It does not assert uniqueness or cuspidality.

The extract does not state levels or congruence subgroups, eigenform hypotheses explicitly, local normalization, good-prime scope, or whether extra cusp, parity or lift-space assumptions are intended. These must be recovered to make a fully specified functorial-lift theorem; no full-level, cuspidal or additional twisting condition is silently imposed.
\subsection{Short English statement}
Seek one genus-$4m$, weight-$k$ Siegel form whose local parameters concatenate those of the two genus-$2m$ forms and multiply their zeroth parameters. The level, eigenform and normalization conventions are missing from the extract.
\subsection{Sources}
\begin{itemize}
\item Supplied report: \emph{Modulformen} (2008), Oberwolfach Reports, DOI 10.4171/owr/2007/51, pp. 3013--3072; proposed year 2007. \url{https://doi.org/10.4171/owr/2007/51}.
\item Local UnsolvedMath registry, permanent ID 30000864, catalogue entry; created 25 August 2026, updated 16 September 2026. Its corrected\_verified note says the lift conjecture and corrupted Satake ranges were restored. The supplied 21 August 2026 literature triage labels the extracted compatibility part of a conjectural functoriality statement, with no located general theorem for the stated correspondence. No additional external reference is supplied; this dated assessment is not independently updated here.
\end{itemize}
\end{document}
